\section{Agent 15-F: primitive recognition modulo radical, attack--heal}
\label{sec:agent15f-primitive-recognition-mod-radical}

\providecommand{\Deltafive}{\Delta_5}
\providecommand{\gDelta}{\mathfrak g_{\Delta_5}}
\providecommand{\Aden}{\mathcal A_{\mathrm{den}}}
\providecommand{\Den}{\mathsf{Den}}
\providecommand{\Prim}{\operatorname{Prim}}
\providecommand{\Primprot}{\operatorname{Prim}_{\mathrm{prot}}}
\providecommand{\Rad}{\operatorname{Rad}}
\providecommand{\sdim}{\operatorname{sdim}}
\providecommand{\smult}{\operatorname{smult}}
\providecommand{\CoHA}{\operatorname{CoHA}}
\providecommand{\Hall}{\mathcal H}
\providecommand{\Ran}{\operatorname{Ran}}
\providecommand{\Spch}{\mathrm{Sp}^{\mathrm{ch}}}
\providecommand{\PhiFA}{\Phi^{\mathrm{FA}}}
\providecommand{\epsor}{\epsilon_o}
\providecommand{\nuDelta}{\nu_{\Delta_5}}

\subsection{Claim attacked}

The attacked theorem is the strongest compact-recognition shortcut:
\[
   \hbox{compact Igusa determinant and signed primitive dimensions}
   \quad\Longrightarrow\quad
   H^0\Primprot(\mathcal F_X)/\Rad \cong \gDelta ,
   \qquad X=K3\times E.
\]
In an even stronger false form, the shortcut says that the surrounding
volumes already supply the compact primitive Lie superalgebra and the
radical quotient because they contain Hall--Drinfeld, bar--cobar, BCOV,
and Borcherds technology.

The attack succeeds.  The determinant, the OP scalar square, the
Gritsenko--Nikulin denominator identity, the formal current envelope, and
the chiral bar--cobar inversion theorem do not construct the compact
primitive bracket.  They also do not determine full parity dimensions,
simple primitive representatives, Chevalley--Serre--orthogonality
relations, a Hopf pairing, or the radical quotient.  Equality of the
denominator is therefore only a necessary numerical shadow.  It is not a
Lie-superalgebra recognition theorem.

\subsection{Local anchors}

\begin{enumerate}
\item \verb|main.tex:188--209| states that primitive recognition is
conditional on a compact Igusa datum and is not an existence theorem.
\item \verb|main.tex:906--934| lists the protected Hall integration
criterion: compact Hall/CoHA sector, oriented extension
correspondences, vanishing cycles, HN completions, protected
integration, primitive projection, and lifted Weyl action.
\item \verb|main.tex:1759--1864| defines the target denominator algebra
\(\Aden=\gDelta\), its root grading, signed root supermultipities, and
denominator \(64^{-1}\Deltafive(2Z)\).
\item \verb|main.tex:1911--1960| makes the \(E_3\)-source and
\(K3\)-to-\(E\) specialization supplied data, then asks for
\[
H^0\Primprot(\mathcal F_X)\stackrel{?}{\cong}\gDelta,\qquad
\Omega_E^{\mathrm{ch}}C_X\stackrel{?}{\simeq}\Den_{\Delta_5,E}.
\]
\item \verb|main.tex:2017--2089| proves the positive-elliptic-degree
locality obstruction and names the two repair paths: projection-finite
locality plus global modes, or a hybrid local/global wrapped
factorization theory.
\item \verb|main.tex:2182--2266| defines the compact Igusa datum
\(\mathfrak K_X\); its \((\Pi)\)-entry is exactly the primitive
recognition certificate.
\item \verb|main.tex:2344--2430| proves only a conditional primitive
recognition theorem: if the presentation-level data, parity dimensions,
generation, and radical quotient are supplied, then the compact
primitive quotient is \(\gDelta\).
\item \verb|main.tex:2605--2752| proves that the denominator determines
signed exponents after the GN root datum is fixed, but forgets
zero-superdimension root spaces, parity dimensions, structure constants,
simple representatives, relation constants, and Hopf-pairing radical.
\item \verb|main.tex:2833--2882| fixes the coefficient dictionary:
\(f(nm,l)\) has three readings, but none is an equality of Hilbert
spaces, brackets, or Hall constants.
\end{enumerate}

The preceding local reports used here are
\[
\begin{gathered}
\verb|agent_material/14_compact_E3_source_from_volumes_attack_heal.tex|,\\
\verb|agent_material/14_ranE_hybrid_wrapping_from_volumes_attack_heal.tex|,\\
\verb|agent_material/14_orientations_protected_integration_from_volumes_attack_heal.tex|,\\
\verb|agent_material/14_primitive_borcherds_recognition_from_volumes_attack_heal.tex|,\\
\verb|agent_material/14_chiral_bar_cobar_Koszul_realization_attack_heal.tex|,\\
\verb|agent_material/14_boundary_rows_physical_hosts_from_volumes_attack_heal.tex|.
\end{gathered}
\]

\subsection{Cross-volume anchors}

\begin{enumerate}
\item \verb|calabi-yau-quantum-groups/FRONTIER.md:31| states the
two-stage architecture
\[
\Phi_d=\Spch_{\Sigma_{d-1},C}\circ\PhiFA_d,
\]
with Stage 2 a specialization, not an inversion.
\item \verb|calabi-yau-quantum-groups/FRONTIER.md:81| states that
bracket-level BPS \(\cong\) BKM remains open beyond the abelian
Heisenberg--Miki core.
\item \verb|calabi-yau-quantum-groups/FRONTIER.md:98--110| fixes the
Hall--Drinfeld double language and says the plain K3 Yangian language is
retracted; it also warns that BKM non-abelianity appears only after the
vertex-operator closure lane.
\item \verb|calabi-yau-quantum-groups/FRONTIER.md:197--205| keeps the
Hall--Drinfeld/BKM-side object as a genuine open frontier: real-root
completion and comparison data remain open.
\item \verb|calabi-yau-quantum-groups/chapters/examples/k3e_bkm_chapter.tex:15--74|
marks the compact Hall realization and completed double as conditional;
the theorem-grade part is the reduced DT character and the
Gritsenko--Nikulin--Borcherds target.
\item \verb|calabi-yau-quantum-groups/chapters/examples/k3e_bkm_chapter.tex:733--753|
records the GN target simple-root and parity convention, including real
even roots and imaginary roots.
\item \verb|calabi-yau-quantum-groups/chapters/examples/k3e_bkm_chapter.tex:970--1022|
records the denominator identity and explicitly separates the primitive
denominator statement from the bar Euler/product bridge problem.
\item \verb|calabi-yau-quantum-groups/notes/compact_k3e_hall_construction_package_20260424.md:63--74|
says the compact Hall--Drinfeld double exists only after the negative
half, Cartan completion, continuous Hopf pairing, radical quotient,
bracket comparison, and centre compatibility are supplied.
\item \verb|calabi-yau-quantum-groups/notes/compact_k3e_hall_construction_package_20260424.md:97--113|
states the same obstruction vector: the positive cone and denominator
normalization do not construct the double.
\item \verb|calabi-yau-quantum-groups/notes/bps_positive_geometry_total_resolution_20260424/agent_06_k3xe_hall_bkm_autborch.md:315--346|
sets a positive-half target whose character follows only after a
PBW-level isomorphism is supplied.
\item \verb|calabi-yau-quantum-groups/notes/bps_positive_geometry_total_resolution_20260424/agent_06_k3xe_hall_bkm_autborch.md:407--451|
keeps positive half separate from Cartan, negative half, pairing,
radical, and the full double.
\item \verb|calabi-yau-quantum-groups/notes/bps_positive_geometry_total_resolution_20260424/agent_06_k3xe_hall_bkm_autborch.md:840--883|
lists finite-height numerical tests and immediately says they do not
prove the bracket.
\item \verb|calabi-yau-quantum-groups/notes/bps_positive_geometry_total_resolution_20260424/agent_06_k3xe_hall_bkm_autborch.md:885--912|
isolates the surviving bracket obstruction \(o_{HB}^{\mathrm{bracket}}\).
\item \verb|chiral-bar-cobar/CLAUDE.md:56--63| and
\verb|chiral-bar-cobar-vol2/CLAUDE.md:348--383| state the
CY-to-chiral direction: Stage 2 is specialization from the CY
factorization object to a curve, never inversion.
\item \verb|chiral-bar-cobar/CLAUDE.md:293--300| forbids conflating
\(A\), \(B(A)\), \(A^!\), and the derived centre; \(\Omega B(A)=A\) is
inversion, not a compact source construction.
\item \verb|chiral-bar-cobar/worldview_synthesis_2026_04_17.tex:239--259|
gives the bar--cobar adjunction and inversion theorem on its proper
domain.  It starts after a curve-level chiral algebra or coalgebra is
given.
\item \verb|chiral-bar-cobar-vol2/FRONTIER.md:58--68| imports the
Hall--Drinfeld double presentation but retains the abelian-at-Lie
discipline.
\item \verb|topological-strings/CLAUDE.md:20--22| says the BCOV volume
is motivation and convention-checking, not a compact CY\(_3\)/BKM
theorem without a matched proof.
\item \verb|topological-strings/main.tex:1194--1250| gives an analogous
source-convention obstruction: a Hamiltonian/PBW source is not the same
as a closed CE-cochain source.
\item \verb|topological-strings/main.tex:2426--2440| states that
Costello's BV formalism supplies a general construction package, while
the paper must still identify the Hamiltonian input and coefficients in
the desired specialization.
\end{enumerate}

\subsection{Attack on the six compact-construction gates}

\paragraph{Gate 1: compact finite-first holomorphic \(E_3\) source.}
Without a supplied charge-completed holomorphic \(E_3\)-factorization
algebra \(\mathcal A^{E_3}_X\), the object
\(\mathcal F_X=\Spch_{K3,E}(\mathcal A^{E_3}_X)\) is only a symbol.
Then \(H^0\Primprot(\mathcal F_X)\) has no source category, no
coproduct, no augmentation, no protected differential, and no completed
topology.  A theorem about its primitives is therefore ill-typed.  The
Vol III \(E_3\) and hCS language is architectural input, not a compact
source construction with finite-first HN quotients in this paper.

\paragraph{Gate 2: Ran(\(E\)) elliptic-degree obstruction.}
Positive elliptic degree dominates \(E\), so ordinary finite-support
\(\Ran(E)\) locality cannot see the full \(s\)-direction.  For primitive
recognition this is more serious than losing a product variable: if the
repair is a decoupled global Fock factor, then it may reproduce the
determinant while giving zero or wrong mixed brackets between vertical
projection-finite primitives and wrapped primitives.  The hybrid repair
must therefore construct mixed extension stacks and prove their
associativity and primitive brackets.  A scalar Fock mode is not a
Borcherds bracket.

\paragraph{Gate 3: oriented compact Hall/CoHA atlas.}
The primitive functor is only defined after a Hall or factorization
object has multiplication, comultiplication, unit, counit, and a
completion compatible with charge.  The reduced \(E\)-quotient cannot be
taken before extension correspondences are built: quotienting too early
erases relative \(E\)-position and can destroy the convolution source of
the bracket.  The compact atlas must therefore live before reduction,
with finite-type HN truncations, vanishing-cycle coefficients, and
Thom--Sebastiani orientation transport.

\paragraph{Gate 4: protected integration and orientation character.}
Protected integration must be a structure-preserving functor, not merely
a scalar trace.  It must commute with Hall product, coproduct,
wall-crossing, primitive projection, parity, and completion.  Otherwise
it can produce the right signed Euler characteristic while forgetting
the primitive Lie bracket and the Hopf pairing.  Constants such as
\(64\) and \(4096\) are normalizations, not orientation-line monodromy.

\paragraph{Gate 5: \(\epsor=\nuDelta\).}
The automorphic character \(\nuDelta\) is theorem-level data of
\(\Deltafive\).  The geometric character \(\epsor\) is undefined until
an oriented Hall object and a lifted Weyl action exist.  Even after
\(\epsor=\nuDelta\) is proved, that equality fixes the sign character
needed by the denominator comparison; it does not by itself construct
simple primitive representatives, Serre relations, imaginary-root
parities, or the Hopf radical quotient.

\paragraph{Gate 6: primitive recognition modulo radical.}
This is the primary obstruction.  Let
\[
   P_X := H^0\Primprot(\mathcal F_X).
\]
The desired object is not \(P_X\) as a graded vector space, but the
radical quotient
\[
   P_X^{\mathrm{red}}
   :=
   P_X/\Rad\langle-,-\rangle
\]
as a charge-graded Lie superalgebra with Cartan, positive and negative
halves, invariant Hopf pairing, and completed topology.  Its
recognition requires a presentation theorem, not a character theorem.

\subsection{First-principles failure modes for gate 6}

\begin{enumerate}
\item \textbf{Grothendieck shadow failure.}  The determinant lives in
the Grothendieck group of graded super vector spaces:
\[
   \prod_\gamma (1-x^\gamma)^{\sdim P_\gamma}.
\]
It is invariant under adding an even and an odd cancelling pair in the
same charge.  Such a pair changes the state space, the possible
brackets, the primitive functor, and the radical, while leaving the
determinant unchanged.

\item \textbf{Zero-bracket failure.}  Give a graded super vector space
the correct \(\sdim P_\gamma=f(nm,l)\) and put the zero bracket on all
positive charges.  Its Fock determinant is still the Igusa product, but
its Lie superalgebra is not \(\gDelta\).

\item \textbf{Wrong-parity failure.}  A signed multiplicity \(a-b\)
does not determine \((a,b)\).  If an imaginary root has signed
multiplicity \(-64\), the determinant cannot distinguish
\((0,64)\) from \((1,65)\).  BKM recognition needs the GN parity split,
not only the superdimension.

\item \textbf{Primitive-versus-PBW failure.}  A product coefficient can
come from a primitive generator or from a PBW monomial in lower
generators after logarithmic/Mobius inversion on rays.  The compact Hall
object must identify which classes are simple primitives and which are
composites.

\item \textbf{Serre-kernel failure.}  A map from the abstract
Borcherds--Kac presentation to \(P_X\) can satisfy the displayed
relations on generators but still have an additional kernel in higher
root degree.  Root superdimensions alone cannot exclude a hidden
zero-superdimension kernel.

\item \textbf{Radical-mismatch failure.}  The completed Hopf pairing can
have a larger radical than the GN/Kac radical, killing target root
spaces, or a smaller radical, leaving extra central primitives.  Either
case gives the same low-order determinant if the additional classes have
zero signed dimension.

\item \textbf{Completion failure.}  Finite-height truncations may have
nondegenerate pairings and PBW bases while the inverse limit develops
\(\varprojlim^1\) cohomology, nonclosed radicals, or discontinuous
brackets.  Recognition is a statement in the completed category.

\item \textbf{Source-convention failure.}  The topological-strings CE
obstruction is the same structural warning in another fabric: a
Hamiltonian/PBW source and a CE-cochain source are not interchangeable.
For Igusa, the Hall commutator source must be fixed before comparing it
with the Borcherds bracket.
\end{enumerate}

\subsection{Exact proof obligations for primitive recognition}

The following obligations are necessary and, together with gates 1--5,
sufficient for the desired recognition statement.

\begin{enumerate}
\item \textbf{Primitive functor.}  Construct a completed, augmented,
charge-graded Hopf or factorization-Hopf object \(\mathcal F_X\) and
define
\[
\Prim(\mathcal F_X)_\gamma
=
\ker\left(
\Delta_\gamma-1\otimes(-)-(-)\otimes1:
(\mathcal F_X)_\gamma\to
\prod_{\gamma_1+\gamma_2=\gamma}
(\mathcal F_X)_{\gamma_1}\widehat\otimes
(\mathcal F_X)_{\gamma_2}
\right)
\]
in the protected completed category.

\item \textbf{Bracket source.}  Prove that the supercommutator of the
oriented Hall product, or equivalently the binary chiral collision
residue, preserves protected primitives after projection:
\[
[-,-]_X:
P_{X,\gamma}\otimes P_{X,\eta}\to P_{X,\gamma+\eta}.
\]

\item \textbf{Cartan and degree map.}  Identify
\[
P_{X,0}\cong \Lambda^{2,1}_{II}\otimes\mathbb C,\qquad
\alpha(n,l,m)=2nf_2-lf_3+2mf_{-2},
\]
and prove that \(\alpha\) sends Hall charge addition to GN root
addition and transports the Euler/BPS pairing to the Lorentzian form
with the manuscript's sign convention.

\item \textbf{Simple representatives.}  Exhibit compact primitive
classes
\[
e_i^X,\quad f_i^X,\quad h_i^X
\]
for the three real simple roots and for every imaginary simple-root
fibre in the GN datum, preserving multiplicity fibres and parity.

\item \textbf{Chevalley relations.}  Prove
\[
[h,e_i^X]=(\alpha_i,h)e_i^X,\qquad
[h,f_i^X]=-(\alpha_i,h)f_i^X,\qquad
[e_i^X,f_j^X]=\delta_{ij}h_i^X
\]
with the correct super signs and normalization of \(h_i^X\).

\item \textbf{Real Serre relations.}  For the even real simple roots,
prove
\[
(\operatorname{ad}e_i^X)^{1-a_{ij}}e_j^X=0,\qquad
(\operatorname{ad}f_i^X)^{1-a_{ij}}f_j^X=0
\]
for \(i\ne j\).  In the visible rank-three real subsystem,
\((\delta_i,\delta_i)=2\) and \((\delta_i,\delta_j)=-2\) for
\(i\ne j\), so \(a_{ij}=-2\) and the visible Serre exponent is \(3\).

\item \textbf{Borcherds orthogonality.}  If \((\alpha_i,\alpha_j)=0\)
and at least one of \(\alpha_i,\alpha_j\) is imaginary, prove
\[
[e_i^X,e_j^X]=0,\qquad [f_i^X,f_j^X]=0
\]
with the super sign rule.

\item \textbf{Parity dimensions.}  Prove equality of full parity
dimensions, not only of superdimensions:
\[
\left(\dim P_{X,\alpha,\overline 0},
      \dim P_{X,\alpha,\overline 1}\right)
=
\left(\dim\gDelta_{\alpha,\overline 0},
      \dim\gDelta_{\alpha,\overline 1}\right)
\]
for all simple imaginary roots and then for all roots after PBW and
radical quotient.

\item \textbf{Generation.}  Prove that \(P_X^{\mathrm{red}}\) is
generated by the Cartan and the simple primitive representatives.

\item \textbf{No extra relations.}  Let
\[
\pi:G(\mathscr B_{\Delta_5})\longrightarrow P_X^{\mathrm{red}}
\]
be the map from the universal Borcherds--Kac superalgebra attached to
the GN datum.  Prove that \(\ker\pi\) is exactly the standard
Borcherds--Serre/radical ideal.

\item \textbf{Hopf pairing.}  Construct positive and negative halves and
a continuous Hopf pairing
\[
\langle-,-\rangle:
P_X^+\widehat\otimes P_X^-\to\mathbb C
\]
compatible with coproduct, Cartan, super signs, and completion.

\item \textbf{Radical quotient.}  Prove
\[
\Rad\langle-,-\rangle=\Rad_{\mathrm{GN/Kac}}
\]
degree by degree and after completion.  The quotient must kill no target
root space and leave no extra central primitive.

\item \textbf{PBW compatibility.}  Prove a completed PBW theorem for the
Hall primitive algebra and show that its PBW filtration is transported
to the PBW filtration on \(U(\gDelta)\).

\item \textbf{Chiral compatibility.}  If the chiral statement is also
claimed, prove that the conilpotent chiral coalgebra \(C_X\) and the
cobar map \(\Omega_E^{\mathrm{ch}}C_X\to\Den_{\Delta_5,E}\) use the
same completion, charge grading, primitive bracket, and radical
quotient.
\end{enumerate}

\subsection{Worked tests needed}

The proof should be attacked at finite height before any pro-limit
claim is made.  For a finite truncation \(T=(N,R)\), define
\[
P_{X,\le T}^{\mathrm{red}}
=
P_{X,\le T}/\Rad_{\le T}.
\]
The required tests are:

\begin{enumerate}
\item \textbf{Coefficient normalization.}  Recompute
\[
\phi_{0,1}=(r^{-1}+10+r)
+q(10r^{-2}-64r^{-1}+108-64r+10r^2)+O(q^2)
\]
and check \(f(0,-1)=1\), \(f(0,0)=10\), \(f(1,1)=-64\),
\(f(1,0)=108\).

\item \textbf{Degree and chamber test.}  Check that every charge in the
finite compact sector maps by \(\alpha(n,l,m)=2nf_2-lf_3+2mf_{-2}\) to
the chosen positive GN cone, and that no active root is represented by
two inequivalent compact charges unless that multiplicity is part of the
GN fibre.

\item \textbf{Primitive/composite split.}  Compute the finite coproduct
matrix and identify primitive classes as the kernel of the reduced
coproduct.  Compare this list with the GN simple root fibres, not with
the product exponents alone.

\item \textbf{Low real-root bracket test.}  For the three real simple
roots \(\delta_1,\delta_2,\delta_3\), compute the compact Hall brackets
\([e_i^X,f_j^X]\) and verify Chevalley normalization.  Then compute the
length-three real Serre brackets
\[
(\operatorname{ad}e_i^X)^3e_j^X,\qquad
(\operatorname{ad}f_i^X)^3f_j^X.
\]

\item \textbf{Imaginary orthogonality test.}  Pick the lowest imaginary
simple roots orthogonal to a real or imaginary simple root and verify
the zero bracket directly in the compact Hall bracket.

\item \textbf{Parity test.}  For the charges with coefficients
\[
f(0,-1)=1,\quad f(0,0)=10,\quad f(1,1)=-64,\quad f(1,0)=108,
\]
compute the actual protected even and odd dimensions, not only their
difference, and compare with the GN parity convention.

\item \textbf{Zero-superdimension audit.}  Search finite height for
nonzero primitive pairs with equal even and odd dimensions.  These are
invisible to the determinant but fatal unless killed by the same
radical as in the GN/Kac construction.

\item \textbf{Pairing-radical matrix test.}  Compute the finite
positive-negative Hopf pairing matrix, its radical, and the induced
Cartan form.  Compare its rank with the GN/Kac quotient at the same
root height.

\item \textbf{PBW Hilbert test.}  After quotienting by the finite
radical, compute the PBW supercharacter from the compact primitives and
compare with the truncation of
\[
\prod_{\gamma}(1-x^\gamma)^{-f(nm,l)}.
\]
Passing this test is necessary but not sufficient; it must be paired
with the bracket and radical tests above.

\item \textbf{Continuity test.}  Verify that the finite radicals,
brackets, and pairings form a Mittag--Leffler inverse system and that
the completed radical is closed.
\end{enumerate}

\subsection{Healed conditional theorem}

\begin{theorem}[Primitive recognition modulo radical, conditional]
Let \(X=K3\times E\).  Assume gates \(1\)--\(5\) have been constructed:
a finite-first charge-completed holomorphic \(E_3\)-source, a
projection-finite plus wrapped or hybrid \(\Ran(E)\) specialization, an
oriented compact Hall/CoHA atlas, protected integration compatible with
Hall product and coproduct, and a lifted Weyl action whose orientation
character satisfies \(\epsor=\nuDelta\) on the type-II Weyl group.

Assume furthermore the gate-\(6\) primitive certificate listed above:
the primitive functor, Hall bracket, Cartan and degree map, simple
representatives, Chevalley relations, real Serre relations, Borcherds
orthogonality, full parity dimensions, generation, no-extra-relations
statement, continuous Hopf pairing, GN/Kac radical quotient, completed
PBW theorem, and completion compatibility.

Then the induced map from the Gritsenko--Nikulin/Borcherds--Kac
presentation is an isomorphism
\[
   H^0\Primprot(\mathcal F_X)/\Rad\langle-,-\rangle
   \cong
   \Aden
   =
   \gDelta
\]
of completed charge-graded Lie superalgebras.  Consequently
\[
   \operatorname{den}
   \bigl(H^0\Primprot(\mathcal F_X)/\Rad\bigr)
   =
   64^{-1}\Deltafive(2Z).
\]
If, in addition, the chiral coalgebra \(C_X\) and the Koszul comparison
\(\Theta_{\mathrm{Kos}}\) are supplied in the same completed category,
then
\[
   \Omega_E^{\mathrm{ch}} C_X
   \simeq
   \Den_{\Delta_5,E}.
\]
\end{theorem}

\begin{proof}[Proof skeleton]
The compact simple representatives define a homomorphism from the
universal Borcherds--Kac Lie superalgebra
\[
G(\mathscr B_{\Delta_5})\longrightarrow P_X^{\mathrm{red}}.
\]
The Cartan, Chevalley, real Serre, orthogonality, and super sign
relations put the Borcherds--Serre ideal in the kernel.  Generation
gives surjectivity after quotient.  The no-extra-relations theorem and
the equality of Hopf radicals identify the kernel exactly with the
standard GN/Kac radical.  Equality of full parity dimensions rules out
invisible zero-superdimension root spaces.  PBW compatibility and
closedness of the completed radical promote the finite-height
identification to the completed algebra.  The denominator statement is
then the existing GN denominator identity for the target algebra, not
the input from which the isomorphism was inferred.
\end{proof}

\subsection{Decision and claim-status recommendation}

The paper should keep
\[
H^0\Primprot(\mathcal F_X)/\Rad \cong \gDelta
\]
as a conditional recognition theorem.  It may state the theorem above
because the assumptions are exactly the missing construction data.  It
should not state or imply that the compact primitive algebra has been
constructed.

The right programme statement is:
\[
\hbox{construct compact primitives and prove the presentation}
\quad\Longrightarrow\quad
\hbox{recover the GN/Borcherds denominator algebra}.
\]
The wrong statement is:
\[
\hbox{recover the Igusa product}
\quad\Longrightarrow\quad
\hbox{compact primitives are the GN/Borcherds algebra}.
\]

\subsection{Open questions}

\begin{enumerate}
\item What is the exact compact category or stability sector whose
objects give \(\mathcal F_X\) before quotienting by \(E\)?
\item Which hybrid local/global extension stacks carry mixed brackets
between projection-finite and wrapped elliptic-degree sectors?
\item Can the strong orientation be made compatible with the Hall
coproduct, not only with multiplication?
\item What is the compact definition of the protected parity split for
imaginary simple roots, including possible zero-superdimension pairs?
\item Does the Hall commutator satisfy the length-three real Serre
relations for \((\delta_i,\delta_j)=-2\) in the visible rank-three
subsystem?
\item Is the Hopf radical closed in the HN/charge completion, and does
it agree degreewise with the GN/Kac radical?
\item Does the inverse system of finite radical quotients have vanishing
\(\varprojlim^1\) obstruction?
\item Is the chiral cobar comparison using the same radical quotient as
the Hall primitive comparison, or only the same denominator character?
\end{enumerate}

\subsection{Files changed and checks}

Files changed by this agent:
\[
\verb|agent_material/15_primitive_recognition_mod_radical_attack_heal.tex|.
\]
No manuscript source, bibliography, generated PDF, branch pointer, or
other agent report was edited.  Checks run before writing: targeted
\verb|rg| and \verb|sed|/line-number reads of \verb|CLAUDE.md|,
\verb|AGENTS.md|, \verb|main.tex|, \verb|proj.bib|, the six
\verb|agent_material/14_*| reports, and the cross-volume anchors listed
above.  Post-write checks run: \verb|git diff --check --| on this file,
targeted readback with \verb|sed|, keyword coverage with \verb|rg|, line
count with \verb|wc -l|, and an ASCII audit.  No global build was run,
since the report is standalone and is not included from \verb|main.tex|.
